\documentclass[11pt]{article}
\usepackage[margin=1in]{geometry}
\usepackage{amsmath,amssymb,amsthm}
\usepackage{booktabs}
\usepackage{hyperref}

\newtheorem{conj}{Conjecture}
\newtheorem{prop}{Proposition}
\theoremstyle{remark}
\newtheorem{remark}{Remark}

\title{Disproof of TLMC Conjecture 00000000462:\\
Spanning Trees of the Circulant Graph $C_n(1,2)$}
\author{TLMC submission}
\date{October 2, 2026}

\begin{document}
\maketitle

\begin{abstract}
TLMC Conjecture 00000000462 claims that $\tau(C_n(1,2))$, the number of spanning
trees of the circulant graph on $n$ vertices with jumps $1$ and $2$, satisfies a
linear recurrence whose largest real characteristic root tends to
$\alpha^2$ with $\alpha = 2+\sqrt{3}$, i.e.\ $(2+\sqrt3)^2 = 7+4\sqrt3 \approx 13.93$.
We evaluate the sequence exactly,
\[
\tau(C_n(1,2)) \;=\; \frac{n\bigl(L_{2n} - 2(-1)^n\bigr)}{5} \qquad (n \ge 5),
\]
with $L_m$ the Lucas numbers, certified against the matrix--tree theorem by exact
integer determinants. Consequently $\tau^{1/n} \to \varphi^2 = (3+\sqrt5)/2 \approx
2.6180$, and every Laplacian eigenvalue is at most $25/4 = 6.25 < 8 < 7+4\sqrt3$.
The claimed limit $\alpha^2 \approx 13.93$ is impossible; the conjecture is \textbf{false}.
\end{abstract}

\section{The conjecture}

\begin{conj}[00000000462]
$\tau(C_n(1,2))$ satisfies a linear recurrence whose largest real characteristic
root tends to $\alpha^2$, where $\alpha = 2+\sqrt{3}$ (``directly verifiable
against the matrix--tree theorem, then stateable in closed form'').
\end{conj}

Throughout, $C_n(1,2)$ is the circulant graph on vertices $\mathbb{Z}_n$ in which
$i$ is adjacent to $i\pm1$ and $i\pm2$; this is a simple graph exactly for $n\ge5$,
the regime relevant to the conjecture.

\section{Exact evaluation}

For $n \ge 5$ the Laplacian eigenvalues of $C_n(1,2)$ are
\[
\lambda_k = 2\bigl(1-\cos\tfrac{2\pi k}{n}\bigr) + 2\bigl(1-\cos\tfrac{4\pi k}{n}\bigr)
= 2(1-c_k)(3+2c_k), \qquad c_k = \cos\tfrac{2\pi k}{n},\ k = 1,\dots,n-1 .
\]
Using the product formula $\prod_{k=0}^{n-1}(x - c_k) = 2^{1-n}\bigl(T_n(x)-1\bigr)$
for the Chebyshev polynomial $T_n$, and $T_n(-\tfrac32) = \tfrac{(-1)^n}{2}L_{2n}$
(because $(-\tfrac32 \pm \tfrac{\sqrt5}{2})^n = (-\varphi^{\mp 2})^n$ and
$\varphi^{2n} + \varphi^{-2n} = L_{2n}$), the matrix--tree theorem yields
\begin{prop}\label{prop:closed}
For all $n \ge 5$,
\[
\tau(C_n(1,2)) = \frac{1}{n}\prod_{k=1}^{n-1}\lambda_k
= \frac{1}{n}\cdot n^2 \cdot \frac{L_{2n} - 2(-1)^n}{5}
= \frac{n\bigl(L_{2n} - 2(-1)^n\bigr)}{5}.
\]
\end{prop}

\begin{remark}[Certification]
Proposition~\ref{prop:closed} was certified computationally, independently of the
Chebyshev algebra: exact integer Bareiss determinants of the Laplacian minor
reproduce the closed form for $n = 5,6,7,8,9,10,12,15,20,24$, and eigenvalue
products match it for $n = 40,80$ with relative error below $10^{-14}$
(script \texttt{reproduce.py} in the submission).
\end{remark}

\begin{table}[h]
\centering
\begin{tabular}{rrrr}
\toprule
$n$ & $\tau(C_n(1,2))$ & $\tau^{1/n}$ & \\
\midrule
5  & $125$ & $2.626528$ & \\
8  & $3{,}528$ & $2.776138$ & \\
15 & $5{,}581{,}500$ & $2.816978$ & (peak region) \\
20 & $915{,}304{,}500$ & $2.805939$ & (verifier: $2.81$) \\
40 & $418{,}891{,}171{,}182{,}561{,}000$ & $2.757735$ & (verifier: $2.76$) \\
80 & $4.3867\times 10^{34}$ & $2.710359$ & (verifier: $2.71$) \\
\bottomrule
\end{tabular}
\caption{Exact values from Proposition~\ref{prop:closed}; the growth rate is
already decreasing toward $\varphi^2 \approx 2.6180$.}
\end{table}

\section{The true growth rate}

Since $L_m = \varphi^m + \psi^m$ with $\psi = -\varphi^{-1}$, $|\psi|<1$, we have
$\varphi^{2n} \le L_{2n} \le \varphi^{2n} + 1$, so Proposition~\ref{prop:closed}
gives the squeeze
\[
\varphi^2 \Bigl(\tfrac{n}{5}\Bigr)^{1/n}
\;\le\; \tau(C_n(1,2))^{1/n} \;\le\;
\varphi^2 \Bigl(1+\varphi^{-2n}\Bigr)^{1/n}\Bigl(\tfrac{n}{5}\Bigr)^{1/n},
\]
hence
\begin{equation}\label{eq:limit}
\lim_{n\to\infty} \tau(C_n(1,2))^{1/n} = \varphi^2 = \frac{3+\sqrt5}{2} \approx 2.618034 .
\end{equation}
Moreover $n \mapsto \tau(C_n(1,2))/n$ is a C-finite sequence with the Lucas
recurrence, whose only characteristic roots are $\varphi^2$ and $-\varphi^{-2}$;
thus the largest real characteristic root of the natural recurrence is
$\varphi^2$ for \emph{every} $n$, not merely in the limit.

\section{Why $\alpha^2$ is impossible}

Each factor $2(1-\cos\theta) \in [0,4]$, so every Laplacian eigenvalue satisfies
$\lambda_k \le 8$; in fact the sharp bound is
\[
\sup_{\theta}\bigl(4 - 2\cos\theta - 2\cos 2\theta\bigr) = \frac{25}{4} = 6.25,
\qquad \text{attained at } \cos\theta = -\tfrac14 .
\]
The matrix--tree theorem therefore gives, for every $n \ge 5$,
\[
\tau(C_n(1,2)) \;\le\; \frac{1}{n}\Bigl(\frac{25}{4}\Bigr)^{n-1}
\;\Longrightarrow\;
\tau(C_n(1,2))^{1/n} < \frac{25}{4} = 6.25 < 8 < \alpha^2 ,
\]
where $\alpha^2 = 7 + 4\sqrt3 > 8$ because $\sqrt3 > \tfrac14$. This contradicts
Conjecture~00000000462 under its stated reading: a limit
$\tau^{1/n} \to \alpha^2 \approx 13.93$ ``directly verifiable against the
matrix--tree theorem'' cannot hold when $\tau^{1/n} < 6.25$ for all $n$ and
$\tau^{1/n} \to \varphi^2 \approx 2.618$ by \eqref{eq:limit}. The separator is
rigorous at the level of integers:
\[
\varphi^2 < 3 < \tfrac{25}{4} < 8 < 9 < \alpha^2 < 14,
\]
using $\sqrt5 < 3$ (as $5 < 9$) and $\tfrac14 < \sqrt3 < \tfrac74$ (as
$1 < 48 < 49 = (4\sqrt3)^2$'s comparison against $1^2$ and $7^2$).

All decisive numerical links are machine-checked in core Lean 4 (no Mathlib,
no axioms): the exact values of $\tau$ for $n \in \{5,\dots,80\}$, the sandwich
$5^{80} < \tau(80)\, 2^{80} \wedge \tau(80) < 8^{80}$ (growth rate in
$(5/2,\,8)$ at $n=80$), the ceiling $a,b \le 4 \Rightarrow a+b \le 8$, and the
integer skeletons $1 < 48 < 49$, $4 < 5$, $5 < 9$. The audit
(\texttt{lake env lean Check.lean}) reports \emph{``does not depend on any
axioms''} for all 25 declarations; there are no \texttt{sorry}s.

\section{Boundary cases and remarks}

\begin{itemize}
\item The closed form of Proposition~\ref{prop:closed} requires $n \ge 5$: for
$n \le 4$ the jumps $1,2,-1,-2$ collide mod $n$. E.g.\ at $n=4$ the graph is
$K_4$ with $\tau = 16$, while the formula would give $36$; at $n=3$ one obtains
a multigraph. These degeneracies are irrelevant to the asymptotic conjecture.
\item The peak of $\tau^{1/n}$ near $n \approx 15$ explains why a naive fit of
early data could suggest a rate well above $\varphi^2$; it remains below $2.82$,
still far from $13.93$.
\item The conjecture's failure is structural: the tree entropy of a
one-dimensional circulant is an integral average of $\log \lambda(\theta)$, which
cannot exceed $\log\max_\theta \lambda(\theta) = \log 6.25 < \log \alpha^2$.
\end{itemize}

\section{Verdict}

\begin{center}
\textbf{Conjecture 00000000462 is FALSE.} The true exponential growth rate is
$\varphi^2 = (3+\sqrt5)/2 \approx 2.618$, with exact closed form
$\tau(C_n(1,2)) = n(L_{2n}-2(-1)^n)/5$ for $n \ge 5$; the claimed root limit
$(2+\sqrt3)^2 \approx 13.93$ exceeds the universal eigenvalue ceiling $25/4$.
\end{center}

\end{document}
